% x86-64 GAS Assembly Quick Reference Card
% Compile: pdflatex card.tex
\documentclass[landscape,a4paper]{article}
\usepackage[landscape,margin=0.3cm,top=0.45cm,bottom=0.3cm]{geometry}
\usepackage{multicol}
\usepackage{booktabs}
\usepackage{array}
\usepackage{xcolor}
\usepackage{tcolorbox}
\usepackage{listings}
\usepackage[T1]{fontenc}
\usepackage{lmodern}

\pagestyle{empty}
\tcbuselibrary{skins}

%%% Compact font sizes %%%
\makeatletter
\renewcommand\normalsize{%
  \@setfontsize\normalsize{7.5pt}{9pt}%
  \abovedisplayskip 4\p@ \@plus1\p@ \@minus1\p@
  \abovedisplayshortskip \abovedisplayskip
  \belowdisplayshortskip \abovedisplayskip
  \belowdisplayskip \abovedisplayskip
}
\renewcommand\small{\@setfontsize\small{7pt}{8.5pt}}
\renewcommand\footnotesize{\@setfontsize\footnotesize{6.5pt}{8pt}}
\renewcommand\scriptsize{\@setfontsize\scriptsize{5.5pt}{6.5pt}}
\renewcommand\tiny{\@setfontsize\tiny{5pt}{6pt}}
\makeatother

%%% Operand type colours %%%
\definecolor{tL}{RGB}{0,80,180}
\definecolor{tM}{RGB}{180,30,30}
\definecolor{tR}{RGB}{0,110,50}
\newcommand{\Lc}{\textcolor{tL}{\textbf{L}}}
\newcommand{\Mc}{\textcolor{tM}{\textbf{M}}}
\newcommand{\Rc}{\textcolor{tR}{\textbf{R}}}

%%% tcolorbox style %%%
\tcbset{
  mysec/.style={
    colback=white, colframe=black!60,
    boxrule=0.2pt, arc=1pt,
    fonttitle=\footnotesize\bfseries,
    left=2pt, right=2pt, top=0pt, bottom=0pt,
    before skip=1pt, after skip=0pt,
  }
}

%%% Listing style %%%
\lstdefinestyle{asmcode}{
  basicstyle=\ttfamily\scriptsize,
  commentstyle=\color{gray}\itshape,
  morecomment=[l]{\#},
  showstringspaces=false,
  breaklines=true,
  frame=none,
  xleftmargin=0pt, xrightmargin=0pt,
  aboveskip=0pt, belowskip=0pt,
  columns=fullflexible,
  keepspaces=true,
  tabsize=4,
}

\begin{document}
\tiny
\setlength{\tabcolsep}{1.5pt}
\setlength{\extrarowheight}{0.5pt}
\setlength{\columnsep}{3pt}
\setlength{\columnseprule}{0.2pt}
\renewcommand{\arraystretch}{0.92}

% ====================== TITLE ======================
\noindent
{\footnotesize\bfseries x86-64 GAS Assembly Quick Reference Card}
\hfill
{\scriptsize
  \Lc\ = Immediate/Literal \quad
  \Mc\ = Memory \quad
  \Rc\ = Register \quad
  \{S\} = Size suffix: b(byte) w(word) l(dword) q(qword)}

\vspace{1pt}

% ===================== PART 1 =====================
\begin{tcolorbox}[mysec,
  title={Part 1: Instruction Set --- AT\&T: \texttt{op src, dst}
         \quad$|$\quad Intel: \texttt{op dst, src}}]

\begin{multicols}{2}
\setlength{\parskip}{0pt}
\setlength{\parsep}{0pt}

% ---------- AT&T Syntax ----------
{\scriptsize\bfseries\sffamily AT\&T Syntax}\par\vspace{-2pt}
{
\begin{tabular}{@{}l l l l p{2.5cm}@{}}
\toprule
\textbf{Mnem.} & \textbf{Op1} & \textbf{Op2} & \textbf{Op3}
  & \textbf{Description} \\
\midrule
\texttt{mov\{S\}}  & \Rc/\Mc/\Lc & \Rc/\Mc & --- & Copy src$\to$dst \\
\texttt{movsx}     & \Rc/\Mc     & \Rc     & --- & Sign-extension move \\
\texttt{movzx}     & \Rc/\Mc     & \Rc     & --- & Zero-extension move \\
\texttt{lea}       & \Mc         & \Rc     & --- & Load effective addr \\
\midrule
\texttt{add\{S\}}  & \Rc/\Mc/\Lc & \Rc/\Mc & --- & $dst \mathrel{+}= src$ \\
\texttt{sub\{S\}}  & \Rc/\Mc/\Lc & \Rc/\Mc & --- & $dst \mathrel{-}= src$ \\
\texttt{imul}      & \Rc/\Mc/\Lc & \Rc     & \Lc
  & Signed mul (1/2/3-op) \\
\texttt{mul\{S\}}  & \Rc/\Mc & --- & --- & Unsigned: RDX:RAX=RAX$*$src \\
\texttt{div\{S\}}  & \Rc/\Mc & --- & --- & Unsigned: RAX=quot,RDX=rem \\
\texttt{idiv\{S\}} & \Rc/\Mc & --- & --- & Signed: RAX=quot,RDX=rem \\
\texttt{inc\{S\}}  & \Rc/\Mc & --- & --- & Increment by 1 \\
\texttt{dec\{S\}}  & \Rc/\Mc & --- & --- & Decrement by 1 \\
\texttt{neg\{S\}}  & \Rc/\Mc & --- & --- & Two's complement negate \\
\midrule
\texttt{and\{S\}}  & \Rc/\Mc/\Lc & \Rc/\Mc & --- & Bitwise AND \\
\texttt{or\{S\}}   & \Rc/\Mc/\Lc & \Rc/\Mc & --- & Bitwise OR \\
\texttt{xor\{S\}}  & \Rc/\Mc/\Lc & \Rc/\Mc & --- & Bitwise XOR \\
\texttt{not\{S\}}  & \Rc/\Mc & --- & --- & Bitwise NOT \\
\midrule
\texttt{shl\{S\}}  & \Lc/\Rc(CL) & \Rc/\Mc & --- & Shift left \\
\texttt{shr\{S\}}  & \Lc/\Rc(CL) & \Rc/\Mc & --- & Shift right (logical) \\
\texttt{sar\{S\}}  & \Lc/\Rc(CL) & \Rc/\Mc & --- & Shift right (arith) \\
\midrule
\texttt{cmp\{S\}}  & \Rc/\Mc/\Lc & \Rc/\Mc & --- & Flags on $dst - src$ \\
\texttt{test\{S\}} & \Rc/\Mc/\Lc & \Rc/\Mc & --- & Flags on $dst\mathbin{\&}src$ \\
\midrule
\texttt{push\{S\}} & \Rc/\Mc/\Lc & --- & --- & Push (RSP$-=$8) \\
\texttt{pop\{S\}}  & \Rc         & --- & --- & Pop (RSP$+=$8) \\
\midrule
\texttt{jmp}       & \Lc/\Mc/\Rc & --- & --- & Unconditional jump \\
\texttt{jcc}       & \Lc         & --- & --- & Conditional jump \\
\texttt{call}      & \Lc/\Mc/\Rc & --- & --- & Call subroutine \\
\texttt{ret}       & ---         & --- & --- & Return \\
\midrule
\texttt{syscall}   & --- & --- & --- & Linux syscall (RAX=\#) \\
\texttt{nop}       & --- & --- & --- & No operation \\
\texttt{cdq/cqo}   & --- & --- & --- & Sign-ext RAX$\to$RDX:RAX \\
\bottomrule
\end{tabular}
}

\columnbreak

% ---------- Intel Syntax ----------
{\scriptsize\bfseries\sffamily Intel Syntax}\par\vspace{-2pt}
{
\begin{tabular}{@{}l l l l p{2.5cm}@{}}
\toprule
\textbf{Mnem.} & \textbf{Op1} & \textbf{Op2} & \textbf{Op3}
  & \textbf{Description} \\
\midrule
\texttt{mov}       & \Rc/\Mc & \Rc/\Mc/\Lc & --- & Copy src$\to$dst \\
\texttt{movsx}     & \Rc     & \Rc/\Mc     & --- & Sign-extension move \\
\texttt{movzx}     & \Rc     & \Rc/\Mc     & --- & Zero-extension move \\
\texttt{lea}       & \Rc     & \Mc         & --- & Load effective addr \\
\midrule
\texttt{add}       & \Rc/\Mc & \Rc/\Mc/\Lc & --- & $dst \mathrel{+}= src$ \\
\texttt{sub}       & \Rc/\Mc & \Rc/\Mc/\Lc & --- & $dst \mathrel{-}= src$ \\
\texttt{imul}      & \Rc     & \Rc/\Mc     & \Lc
  & Signed mul (1/2/3-op) \\
\texttt{mul}       & \Rc/\Mc & --- & --- & Unsigned: RDX:RAX=RAX$*$src \\
\texttt{div}       & \Rc/\Mc & --- & --- & Unsigned: RAX=quot,RDX=rem \\
\texttt{idiv}      & \Rc/\Mc & --- & --- & Signed: RAX=quot,RDX=rem \\
\texttt{inc}       & \Rc/\Mc & --- & --- & Increment by 1 \\
\texttt{dec}       & \Rc/\Mc & --- & --- & Decrement by 1 \\
\texttt{neg}       & \Rc/\Mc & --- & --- & Two's complement negate \\
\midrule
\texttt{and}       & \Rc/\Mc & \Rc/\Mc/\Lc & --- & Bitwise AND \\
\texttt{or}        & \Rc/\Mc & \Rc/\Mc/\Lc & --- & Bitwise OR \\
\texttt{xor}       & \Rc/\Mc & \Rc/\Mc/\Lc & --- & Bitwise XOR \\
\texttt{not}       & \Rc/\Mc & --- & --- & Bitwise NOT \\
\midrule
\texttt{shl}       & \Rc/\Mc & \Lc/\Rc(CL) & --- & Shift left \\
\texttt{shr}       & \Rc/\Mc & \Lc/\Rc(CL) & --- & Shift right (logical) \\
\texttt{sar}       & \Rc/\Mc & \Lc/\Rc(CL) & --- & Shift right (arith) \\
\midrule
\texttt{cmp}       & \Rc/\Mc & \Rc/\Mc/\Lc & --- & Flags on $dst - src$ \\
\texttt{test}      & \Rc/\Mc & \Rc/\Mc/\Lc & --- & Flags on $dst\mathbin{\&}src$ \\
\midrule
\texttt{push}      & \Rc/\Mc/\Lc & --- & --- & Push (RSP$-=$8) \\
\texttt{pop}       & \Rc         & --- & --- & Pop (RSP$+=$8) \\
\midrule
\texttt{jmp}       & \Lc/\Mc/\Rc & --- & --- & Unconditional jump \\
\texttt{jcc}       & \Lc         & --- & --- & Conditional jump \\
\texttt{call}      & \Lc/\Mc/\Rc & --- & --- & Call subroutine \\
\texttt{ret}       & ---         & --- & --- & Return \\
\midrule
\texttt{syscall}   & --- & --- & --- & Linux syscall (RAX=\#) \\
\texttt{nop}       & --- & --- & --- & No operation \\
\texttt{cdq/cqo}   & --- & --- & --- & Sign-ext RAX$\to$RDX:RAX \\
\bottomrule
\end{tabular}
}

\end{multicols}

\vspace{-2pt}
{\scriptsize
\textbf{jcc variants:}\quad
\texttt{je}/\texttt{jz}~($=$)\quad
\texttt{jne}/\texttt{jnz}~($\neq$)\quad
\texttt{jg}~(signed$>$)\quad
\texttt{jl}~(signed$<$)\quad
\texttt{jge}~($\geq$)\quad
\texttt{jle}~($\leq$)\quad
\texttt{ja}~(unsigned$>$)\quad
\texttt{jb}~(unsigned$<$)\quad
\texttt{jae}~($\geq$)\quad
\texttt{jbe}~($\leq$)\quad
\texttt{js}~(SF=1)\quad
\texttt{jns}~(SF=0)}

\end{tcolorbox}

% ===================== PART 2 =====================
\begin{tcolorbox}[mysec,
  title={Part 2: General Purpose Registers --- x86-64}]

{
\begin{tabular}{@{}l l l l l p{3.8cm}@{}}
\toprule
\textbf{8-bit} & \textbf{16-bit} & \textbf{32-bit} & \textbf{64-bit}
  & \textbf{Save} & \textbf{Purpose / Calling Convention} \\
\midrule
AL & AX & EAX & RAX & No  & Accumulator; return value \\
BL & BX & EBX & RBX & Yes & Base; callee-saved \\
CL & CX & ECX & RCX & No  & Counter; Win:arg1/SysV:arg4 \\
DL & DX & EDX & RDX & No  & Data; Win:arg2/SysV:arg3 \\
SIL & SI & ESI & RSI & No & Source idx; SysV:arg2 \\
DIL & DI & EDI & RDI & No & Dest idx; SysV:arg1 \\
SPL & SP & ESP & RSP & --- & Stack pointer \\
BPL & BP & EBP & RBP & Yes & Frame ptr; callee-saved \\
\midrule
R8B  & R8W  & R8D  & R8  & No  & arg5 (Win+SysV) \\
R9B  & R9W  & R9D  & R9  & No  & arg6 (Win+SysV) \\
R10B & R10W & R10D & R10 & No  & caller-saved; Win:static chain \\
R11B & R11W & R11D & R11 & No  & caller-saved; Win:tail-call tmp \\
R12B & R12W & R12D & R12 & Yes & callee-saved \\
R13B & R13W & R13D & R13 & Yes & callee-saved \\
R14B & R14W & R14D & R14 & Yes & callee-saved \\
R15B & R15W & R15D & R15 & Yes & callee-saved \\
\midrule
--- & FLAGS & EFLAGS & RFLAGS & --- & CF(carry) ZF(zero) SF(sign) OF(overflow) PF(parity) \\
--- & IP    & EIP    & RIP    & --- & Instruction pointer \\
\bottomrule
\end{tabular}
}

\vspace{-1pt}
{\scriptsize
\textbf{Win x64:} RCX,RDX,R8,R9$\to$stack (32B shadow)\qquad
\textbf{Linux SysV:} RDI,RSI,RDX,RCX,R8,R9$\to$stack\qquad
\textit{Ret: RAX(64b) / RDX:RAX(128b)}}

\end{tcolorbox}

% ===================== PART 3 =====================
\begin{tcolorbox}[mysec,
  title={Part 3: GAS Program Templates (AT\&T Syntax, x86-64 Linux)}]

\begin{multicols}{2}
\setlength{\parskip}{0pt}

% ---------- Template 1: libc ----------
{\scriptsize\bfseries\sffamily Template 1: With C Library (libc)}\par
\vspace{1pt}
\begin{lstlisting}[style=asmcode]
# Build: gcc tmpl.s -o tmpl
#-------------------------------------------
    .section .data          # [Data Segment]
fmt: .asciz "Result: %d\n"

    .section .bss           # [BSS Segment]
buf: .space 256

    .section .text          # [Text Segment]
    .extern printf          # [Extern Declarations]
    .extern exit

# [Function Segment]
# helper:
#   pushq  %rbp
#   movq   %rsp, %rbp
#   ... body ...
#   popq   %rbp
#   ret

# [Main Entry Point]
    .globl main
# Windows: .globl _main
# Windows: _main:
main:
    pushq  %rbp
    movq   %rsp, %rbp
    leaq   fmt(%rip), %rdi  # arg1 (SysV)
    movl   $42, %esi        # arg2 (SysV)
    # Win: mov rcx, offset fmt
    # Win: mov edx, 42
    xorl   %eax, %eax
    call   printf
    xorl   %eax, %eax       # return 0
    leave
    ret
\end{lstlisting}

\columnbreak

% ---------- Template 2: standalone ----------
{\scriptsize\bfseries\sffamily Template 2: Standalone (syscall)}\par
\vspace{1pt}
\begin{lstlisting}[style=asmcode]
# Build: as --64 tmpl.s -o tmpl.o
#         ld tmpl.o -o tmpl
#-------------------------------------------
    .section .data          # [Data Segment]
msg: .ascii "Hello!\n"
.equ msglen, . - msg

    .section .bss           # [BSS Segment]

    .section .text          # [Text Segment]
    # (no extern for syscall-only)

# [Function Segment]
print_msg:
    pushq  %rbp
    movq   %rsp, %rbp
    movq   $1, %rax         # sys_write
    movq   $1, %rdi         # stdout
    leaq   msg(%rip), %rsi
    movq   $msglen, %rdx
    syscall
    popq   %rbp
    ret

# [Main Entry Point]
    .globl _start
# Windows: .globl mainCRTStartup
# Windows: mainCRTStartup:
_start:
    call   print_msg
    movq   $60, %rax        # sys_exit
    xorq   %rdi, %rdi
    syscall
\end{lstlisting}

\end{multicols}

\vspace{-1pt}
{\scriptsize
\textbf{Key:} \texttt{main}=libc (printf,malloc...);
\texttt{\_start}=bare Linux (syscall only, no libc).
RSP must be 16-byte aligned before \texttt{call}.
\quad\textit{Syscalls:} read=0 write=1 open=2 close=3 exit=60.}

\end{tcolorbox}

\end{document}
